\chapter{Spoken and Colloquial Iridian}\label{ch:spoken}

\section{Introduction}

This chapter sketches a small part of the spoken and urban colloquial register
of present-day Iridian. The written standard remains the reference variety of
this grammar, but it would be misleading to treat everyday spoken usage as
merely a noisier version of that standard. In ordinary conversation, speakers
make frequent use of reduced connectives, clause-peripheral hedges, semantically
bleached nouns, and other high-frequency colloquial forms whose contribution is
partly lexical and partly interactional.

The treatment here is explicitly synchronic and non-exhaustive. Several spoken
phenomena have already been noted elsewhere in the grammar, for example the
contraction of sequential \ird{nesu} to colloquial \ird{nu} in
Chapter~\ref{ch:complex} and the intentional-future use of the purposive
converb in the same chapter. What follows does not attempt a complete grammar
of casual speech. It instead documents a small set of common colloquial
morphemes whose behavior is especially revealing for the interaction of lexical
bleaching, discourse management, and ordinary Iridian morphosyntax.

\section{Common colloquial morphemes}\label{sec:spoken-morphemes}

\subsection{The connective \ird{nu}}\label{sec:spoken-nu}

The colloquial connective \ird{nu} has already been introduced in
Section~\ref{sec:seq-copula} of Chapter~\ref{ch:complex}. Synchronically it is
best treated as the reduced spoken continuation of sequential \ird{nesu}, but
its discourse function is now wider than simple sequential clause-linking. In
rapid speech it often serves as a light connective between clauses, turns, or
discourse steps, with meanings ranging from loose \trsl{and then} to weaker
\trsl{so}, \trsl{well then}, or simple transition management.

What matters for the present chapter is that \ird{nu} belongs to colloquial
discourse management, not to the rigid preverbal particle hierarchy of
Chapter~\ref{ch:function-words}. It sits at the edge of clauses and utterances,
helps the speaker move the discourse along, and shows how readily spoken
Iridian recruits reduced forms into connective work.

\subsection{\ird{Žema} and reduced \ird{žma}}\label{sec:spoken-zema}

Among colloquial noun-like items, \ird{žema} is especially revealing. In
present-day spoken Iridian it is best analysed as a semantically bleached
colloquial lexeme with three closely related functions. First, it may serve as
a placeholder noun meaning roughly \trsl{thingy}, \trsl{that thing}, or
\trsl{what-do-you-call-it}. Secondly, it may function as a dismissive or
evaluative noun meaning \trsl{nonsense}, \trsl{excuses}, or \trsl{all that
fuss}. Thirdly, it may appear at the edge of an utterance as a playful hedge,
roughly \trsl{not really}, \trsl{just kidding}, or \trsl{I am only saying it
lightly}.

Alongside the fuller form \ird{žema}, colloquial speech also shows a more
reduced variant \ird{žma}. The reduced form is faster, more interactional, and
more typical of relaxed or intimate conversation, but it is not a different
lexeme with different core meanings. Rather, it is a spoken reduction of the
same colloquial item. Speakers may therefore choose \ird{žema} in clearer or
less reduced casual speech, while \ird{žma} is especially natural in rapid,
playful, or dismissive turns.

\subsubsection{Category and core values}

The basic grammatical point is that \ird{žema} behaves as a noun-like base in
ordinary Iridian syntax. In its placeholder use it fills a nominal slot without
contributing a stable lexical class of its own:

\pex
\a
\begingl
\gla Tak žema Janku koupnik.//
\glb this thingy Janek.\Trs{} buy-\Trs{}-\Pf{}//
\glft \trsl{Janek bought that thing / that whatchamacallit.}//
\endgl
\a
\begingl
\gla Tak žema.//
\glb this nonsense//
\glft \trsl{That is nonsense / that is just an excuse.}//
\endgl
\xe

In the first example, \ird{žema} is simply a placeholder noun occupying the
patient slot of an ordinary transitive clause. In the second, the same form is
interpreted dismissively: the speaker rejects what has just been said or done as
empty, unserious, or merely performative. The lexical content is therefore weak
in both uses, but the syntactic behavior remains nominal.

\subsubsection{Plurality and expressive reduplication}

The nominal behavior of the item becomes still clearer when it combines with the
plural clitic \ird{ne}. As elsewhere in the noun phrase
(§\,\ref{sec:nm-num}), \ird{ne} does not force a reading of neatly countable
plurality. With \ird{žema}, it usually denotes a loose aggregate: \trsl{all
those things}, \trsl{all that stuff}, or, in a more evaluative context, \trsl{all
that nonsense}.

\pex
\a
\begingl
\gla Ne žema Marek koupnik.//
\glb \Pl{} thingy Marek.\Trs{} buy-\Trs{}-\Pf{}//
\glft \trsl{Marek bought all that stuff.}//
\endgl
\a
\begingl
\gla Ne žemažema Marek koupnik.//
\glb \Pl{} thingy-thingy Marek.\Trs{} buy-\Trs{}-\Pf{}//
\glft \trsl{Marek bought all that random junk / all that nonsense stuff.}//
\endgl
\xe

The reduplicated form \ird{žemažema} adds an expressive nuance of vagueness,
theatricality, or dismissal. It is more overtly colloquial than plain
\ird{žema} and is especially natural when the speaker wishes to downplay the
importance of the referents or treat them as an undifferentiated heap of extras,
details, or irritations. In rapid speech the same pattern may surface as reduced
\ird{žmažma}, but the fuller reduplicated form is the clearest citation shape.

A further compressed form \ird{žežma} appears in particularly rapid or
intimate speech, resulting from syncope of the medial syllable of the full
reduplicated shape: \ird{žemažema} contracts to \ird{žežma}. This mirrors
the reduction of the simplex \ird{žema} to \ird{žma} and has a comparable
sociolinguistic distribution---casual, fast-paced, or interactionally marked
registers---without any change in meaning. The three members \ird{žemažema},
\ird{žmažma}, and \ird{žežma} are therefore best treated as stylistic rather
than semantic variants of the reduplicated base.

Informal writing and colloquial lexicographic listings often record the plural
combinations as single orthographic words: \ird{nežma} (for \ird{ne} +
\ird{žma}) and \ird{nežežma} (for \ird{ne} + \ird{žežma}). These spellings
reflect the phonological cohesion of the plural clitic with the base in fast
speech; grammatically, \ird{ne} retains its ordinary clitic status. The
gradient of expressivity noted above is preserved in the plural: \ird{nežma}
is the more neutral aggregate plural (\trsl{all that stuff / all those things}),
while \ird{nežežma} is more overtly evaluative (\trsl{all that performative
nonsense / all that random fussing}), carrying into the plural the heightened
dismissiveness of the reduplicated base.

\subsubsection{Clause-peripheral hedge use}

The same colloquial base may also appear clause-peripherally, especially in the
reduced form \ird{žma}. In this use it no longer behaves like an ordinary noun
inside the clause, but as a parenthetical hedge by which the speaker playfully
withdraws, softens, or reframes what has just been said. Even so, it should not
be treated as a member of the formal particle hierarchy. Its contribution is
interactional and parenthetical rather than structural in the narrow sense.

\pex
\begingl
\gla Janek v ščenik, žma.//
\glb Janek.\Dir{} already arrive-\Pf{} hedge//
\glft \trsl{Janek has already arrived, not really / just kidding.}//
\endgl
\xe

This utterance-edge use follows naturally from the same semantic bleaching seen
in the nominal uses. The speaker introduces a proposition and then immediately
marks it as playful, non-committal, or not to be taken at full face value.
Synchronically, then, \ird{žema} and \ird{žma} form a small colloquial system:
an underspecified noun in the syntax, an evaluative noun in stance-marking
predication, and a parenthetical hedge in clause-peripheral discourse use.

\subsubsection{Verbal use and productive recategorisation}\label{sec:spoken-zema-verbal}

The most grammatically productive extension of the \ird{žema} system is the
recruitment of the nominal base as a colloquial verbal stem. These verbal uses
are best understood not as a single fixed lexical entry but as an instance of
productive recategorisation: the nominal base is assigned a thematic class and
then set to work with ordinary Iridian valency and aspect-alignment morphology.
This is a familiar process in colloquial Iridian, where emotionally vivid or
semantically bleached bases are frequently pressed into verbal service without
any formal derivational affix mediating the transition.

The colloquial verbal stem is \ird{žemat-}, assigned to thematic class~I
(thematic consonant \ird{-t}), with antipassive stem \ird{žemač-}. The
assignment to class~I reflects a well-attested tendency for recently verbalized
colloquial bases to adopt the most common and unmarked thematic alternation.
As a verbal stem, \ird{žemat-} takes the full complement of valency and
aspect-alignment morphology described in Chapter~\ref{ch:verbs}. What remains
unusual is not the morphological behaviour---it is entirely transparent---but
the semantics: the lexical content of the stem is no more determinate than that
of the nominal base from which it derives.

In its actor-focused (antipassive) use, the stem describes the agent as
engaging in vague, performative, evasive, or unserious activity. The core
meaning is most readily glossed as \trsl{to mess around}, \trsl{to make
excuses}, \trsl{to put on a show}, or \trsl{to fuss about without achieving
anything}, though the exact nuance depends on context in the same way as the
nominal base does. The antipassive perfective \ird{žemačék} reports a completed
episode of such activity; the antipassive retrospective \ird{žemačaní} presents
it with continuing present relevance and a characteristically evaluative
charge---the speaker frames the prior activity as having amounted to nothing or
as having been mere performance; the continuous \ird{žemata} describes it as
an ongoing state or habitual disposition:

\pex
\a
\begingl
\gla Anka žemačék.//
\glb Anka.\Dir{} žemat-\Antip{}-\Pf{}//
\glft \trsl{Anka was just messing around / making excuses.}//
\endgl
\a
\begingl
\gla Anka žemačaní.//
\glb Anka.\Dir{} žemat-\Antip{}-\Ret{}//
\glft \trsl{Anka has been making nothing but excuses.}//
\endgl
\a
\begingl
\gla Anka žemata.//
\glb Anka.\Dir{} žemat-\Cont{}//
\glft \trsl{Anka is always fussing / is the type to make excuses.}//
\endgl
\xe

In its patient-focused (ergative) use, the stem \ird{žemat-} takes the
transitive marker \ird{-n-} and acquires a closely related but distinct
reading: \trsl{to handle carelessly}, \trsl{to give only a vague or makeshift
treatment to}, \trsl{to fudge}, or \trsl{to treat as a trivial matter}. In
the ergative aspects the agent appears in the transitive case and the patient
in the directional:

\pex
\a
\begingl
\gla Anka tak žematnik.//
\glb Anka.\Trs{} that.\Dir{} žemat-\Trs{}-\Pf{}//
\glft \trsl{Anka fudged that / handled that carelessly.}//
\endgl
\a
\begingl
\gla Anka tak žematnaní.//
\glb Anka.\Trs{} that.\Dir{} žemat-\Trs{}-\Ret{}//
\glft \trsl{Anka has been treating that carelessly the whole time.}//
\endgl
\xe

The same verbal suffixes attach to the reduplicated base \ird{žemažema-}
(antipassive stem \ird{žemažemač-}) or to its compressed variant \ird{žežma-}
(antipassive stem \ird{žežmač-}). The reduplicated verbal is semantically
parallel to the reduplicated nominal: it is more overtly expressive and more
dismissive than the simple verbal form, conveying not merely an isolated episode
of vague activity but a sustained, layered, or theatrically amplified pattern
of performative moves:

\pex
\a
\begingl
\gla Anka žemažemačék.//
\glb Anka.\Dir{} \textsc{rdp}-žemat-\Antip{}-\Pf{}//
\glft \trsl{Anka went through all her usual vague fussing.}//
\endgl
\a
\begingl
\gla Anka žežmata.//
\glb Anka.\Dir{} \textsc{rdp}-žemat-\Cont{}//
\glft \trsl{Anka is always going through all that performative nonsense.}//
\endgl
\xe

Semantically, all verbal uses share a stable core even as contextual readings
vary widely. The predicate never names a concrete or specific action; it instead
attributes to the agent a \emph{stance}---vagueness, evasiveness, theatricality,
or dismissal---and that stance is precisely what the nominal base also
contributes. \ird{Žemačék} in one context means \trsl{she was just messing
around}; in another, \trsl{he was making excuses}; in a third, \trsl{they were
fussing over trivial extras}. What remains constant is the evaluative frame:
the speaker characterises the activity not by what was done but by how it was
done, and the \textit{how} is always some variant of unserious, unfocused, or
merely performative.

It should further be noted that the verbal morphology is entirely transparent.
The forms produced from \ird{žemat-} follow all ordinary Iridian verbal rules:
class~I mutation of the thematic consonant under the antipassive
(\ird{-t} → \ird{-č}), ergative-absolutive alignment in the perfective and
retrospective, nominative-accusative alignment in the continuous and
progressive, and the regular fusion of aspect and alignment into a single
suffix. The colloquial character of the item lies entirely in its origin and
semantics; its morphological behaviour is unremarkable. This is the hallmark of
productive recategorisation: a semantically bleached colloquial noun has entered
verbal morphology without requiring any special accommodation from the
language's ordinary morphological machinery.
